\chapter{The IOS Command Line}\label{ch:cli}
\bp{2.4}

\begin{outcomes}
By the end of this chapter you will be able to:
\begin{tight}
  \item \textbf{Move} between IOS command modes without hesitating, and recognise
        which mode a prompt indicates.
  \item \textbf{Distinguish} the running configuration from the startup
        configuration, and save deliberately rather than by habit.
  \item \textbf{Perform} a first-time device setup: hostname, addressing,
        passwords, SSH and a banner.
  \item \textbf{Filter} \cmd{show} output, which is what makes long output usable
        under exam time pressure.
\end{tight}
\end{outcomes}

\begin{prereq}
\chref{ipv4} for addressing and \chref{transport} for the diagnostic commands
you will now be typing for real.
\end{prereq}

\begin{keyterms}
user EXEC $\bullet$ privileged EXEC $\bullet$ global configuration $\bullet$
interface configuration $\bullet$ line configuration $\bullet$ running-config
$\bullet$ startup-config $\bullet$ NVRAM $\bullet$ context-sensitive help
$\bullet$ pipe $\bullet$ banner $\bullet$ SSH
\end{keyterms}

\section{Modes, and the prompt that tells you where you are}

\begin{longtable}{@{}L{4.6cm}L{3.6cm}L{6.8cm}@{}}
\toprule
\rowcolor{primary!15}
\textbf{Mode} & \textbf{Prompt} & \textbf{What you can do, and how you got there} \\
\midrule
\endhead
User EXEC & \prompt{SW1>} & Look, barely. Enter with \cmd{enable} \\
\rowcolor{lightbg}
Privileged EXEC & \prompt{SW1\#} & All \cmd{show} commands, \cmd{copy},
\cmd{reload}, \cmd{debug}. Enter configuration with \cmd{configure terminal} \\
Global configuration & \prompt{SW1(config)\#} & Device-wide settings: hostname,
users, routing \\
\rowcolor{lightbg}
Interface configuration & \prompt{SW1(config-if)\#} & One interface. Entered with
\cmd{interface Gi0/1} \\
Line configuration & \prompt{SW1(config-line)\#} & Console and VTY access.
Entered with \cmd{line vty 0 15} \\
\rowcolor{lightbg}
VLAN configuration & \prompt{SW1(config-vlan)\#} & Naming a VLAN.
\chref{vlans} \\
Router configuration & \prompt{R1(config-router)\#} & A routing protocol.
\chref{ospfv2} \\
\bottomrule
\end{longtable}

\begin{conceptbox}[Getting out, and the two ways to do it]
\cmd{exit} steps back one level. \cmd{end}, or \cmd{Ctrl-Z}, jumps straight to
privileged EXEC from anywhere. When an exam question shows you a prompt, the
prompt is telling you which commands are currently legal --- a global command
typed in interface mode is rejected, and that is a favourite trap.

One shortcut worth knowing: \cmd{do} runs a privileged EXEC command without
leaving configuration mode, so \cmd{do show ip interface brief} works from
\prompt{SW1(config-if)\#}.
\end{conceptbox}

\section{Two configurations, and only one survives a reboot}

\begin{definitionbox}[running-config and startup-config]
The \term{running configuration} is in RAM. It is what the device is doing right
now, and every command you type takes effect there \emph{immediately}.

The \term{startup configuration} is in NVRAM. It is what the device will load
when it next boots. Nothing you type reaches it until you copy it there.
\begin{tight}
  \item \cmd{copy running-config startup-config} --- save. Abbreviated
        \cmd{wr} or \cmd{write memory} on most platforms.
  \item \cmd{show running-config} --- what is in force.
  \item \cmd{show startup-config} --- what will be in force after a reload.
\end{tight}
\end{definitionbox}

\begin{alertbox}[Not saving is a technique, not only a mistake]
Because configuration takes effect immediately and is lost on reload, a
deliberately unsaved change is a safety net. If you are about to modify the
access list or the interface you are connected through, and you get it wrong, a
power cycle restores the last saved configuration and your access with it.

The professional habit is: make the risky change, verify you still have access,
\emph{then} save. Reversing that order is how people lock themselves out of
devices in another building.
\end{alertbox}

\section{Getting help, and typing less}

\begin{longtable}{@{}L{4.6cm}L{11.0cm}@{}}
\toprule
\rowcolor{primary!15}
\textbf{Keystroke} & \textbf{Effect} \\
\midrule
\endhead
\cmd{?} & List every command valid here \\
\rowcolor{lightbg}
\cmd{sh ?} & List commands starting with ``sh'' \\
\cmd{show ip ?} & List the next keyword. Note the space --- \cmd{show ip?} means
something different \\
\rowcolor{lightbg}
\cmd{Tab} & Complete the current keyword \\
Abbreviation & Type the shortest unambiguous form. \cmd{sh ip int br} is
\cmd{show ip interface brief} \\
\rowcolor{lightbg}
\cmd{Ctrl-A} / \cmd{Ctrl-E} & Jump to start / end of line \\
\cmd{Ctrl-C} & Abandon this command \\
\rowcolor{lightbg}
\cmd{Ctrl-Shift-6} & Abort a running ping or traceroute \\
Up arrow & Previous command \\
\bottomrule
\end{longtable}

\begin{conceptbox}[Filtering output is a survival skill]
A full \cmd{show running-config} on a production switch runs to hundreds of
lines. The pipe cuts it to what you need, and in an exam it is the difference
between answering and scrolling.
\begin{tight}
  \item \cmd{| include vlan} --- only lines containing ``vlan''
  \item \cmd{| exclude \textasciicircum!} --- drop the \cmd{!} comment lines
  \item \cmd{| begin interface} --- start output at the first match
  \item \cmd{| section interface Gi0/1} --- that interface's whole block
\end{tight}
\end{conceptbox}

\begin{verify}{SW1}{show running-config | section interface GigabitEthernet0/1}
interface GigabitEthernet0/1
 description Uplink to SW2
 switchport trunk encapsulation dot1q
 switchport mode trunk
 switchport trunk allowed vlan 10,20,30
\end{verify}

\section{A first-time setup, end to end}

\begin{config}{a switch from the box to remotely manageable}
enable
configure terminal
!
! 1. Identify the device. The hostname appears in every prompt and log.
hostname SW1
!
! 2. Protect privileged mode. "secret" is hashed; "password" is not --
!    never use "enable password".
enable secret Str0ng-Enable-Secret
!
! 3. A local user for SSH, with full privilege.
username netadmin privilege 15 secret Str0ng-User-Secret
!
! 4. A management address. A layer 2 switch has no routed ports, so
!    this lives on an SVI, and it needs a gateway to be reachable
!    from another subnet.
interface Vlan1
 ip address 10.0.0.11 255.255.255.0
 no shutdown
exit
ip default-gateway 10.0.0.1
!
! 5. SSH needs a domain name and a key. 2048-bit minimum; version 2 only.
ip domain-name example.local
crypto key generate rsa modulus 2048
ip ssh version 2
!
! 6. Console: require login, and stop output interrupting your typing.
line console 0
 login local
 logging synchronous
 exec-timeout 10 0
exit
!
! 7. Remote access: SSH only, authenticated against local users.
line vty 0 15
 login local
 transport input ssh
 exec-timeout 10 0
exit
!
! 8. A legal banner. Never write "welcome" -- it has been argued in court.
banner motd # Authorised access only. Activity is logged. #
!
end
copy running-config startup-config
\end{config}

\begin{verify}{SW1}{show ip ssh}
SSH Enabled - version 2.0
Authentication timeout: 120 secs; Authentication retries: 3
\end{verify}

\begin{pitfall}
\begin{tight}
  \item \textbf{Using \cmd{enable password}.} It is stored in clear text.
        \cmd{enable secret} is hashed, and where both exist the secret wins.
  \item \textbf{Forgetting \cmd{transport input ssh}.} Leave it at the default
        and Telnet stays available, which undoes the point of configuring SSH.
  \item \textbf{Generating a key without a hostname and domain name.} The key is
        named after them, and the command fails or produces a useless key if they
        are still at defaults.
  \item \textbf{Saving before verifying.} See the warning above. Verify access,
        then save.
  \item \textbf{Omitting \cmd{logging synchronous}.} Without it a log message
        arrives in the middle of the line you are typing, and you cannot tell
        what you have actually entered.
\end{tight}
\end{pitfall}

\begin{breakfix}[The change worked yesterday and the device came back without it.]
A colleague configured a new VLAN and confirmed it working. After a power cut
overnight, the VLAN is gone and the configuration looks like it did last week.

\textbf{Diagnosis.} The change was made in the running configuration and never
copied to the startup configuration. The reload loaded NVRAM, which did not have
it. Nothing is broken --- the device did exactly what it was told.

\textbf{Confirm it} by comparing the two before doing anything else:

\begin{verify}{SW1}{show archive config differences nvram:startup-config system:running-config}
!Contextual Config Diffs:
+vlan 30
+ name Guest
+interface GigabitEthernet0/9
+ switchport access vlan 30
\end{verify}

Any output at all means the two differ and a reload will lose the difference.

\textbf{Fix.} Re-apply the change, verify, then \cmd{copy running-config
startup-config}. As a habit, run that comparison before you leave any device.
\end{breakfix}

\begin{lab}{Bring up a switch you can reach from your desk}{Packet Tracer}
\textbf{Topology.} One switch, one PC on the same VLAN, connected by console
cable and by an Ethernet cable.

\begin{tasks}
  \item Connect by console. Move through all five modes in turn and write down
        the prompt at each. Use \cmd{end} to jump back from the deepest.
  \item Set hostname, \cmd{enable secret} and a local user.
  \item Give \cmd{Vlan1} an address in the PC's subnet and a default gateway.
        Ping the switch from the PC to prove it.
  \item Enable SSH properly: domain name, 2048-bit key, version 2, and
        \cmd{transport input ssh} on the VTY lines.
  \item SSH from the PC. Then try Telnet and confirm it is refused.
  \item Make a small change, do \emph{not} save, reload, and observe the change
        vanish. Repeat with a save.
  \item \textbf{Extension.} Practise \cmd{| include}, \cmd{| section} and
        \cmd{| begin} on \cmd{show running-config} until you can extract a single
        interface block in one command without thinking.
\end{tasks}
\end{lab}

\begin{checkpoint}
\begin{tightnum}
  \item You are at \prompt{SW1(config-if)\#} and need to check the routing table
        without losing your place. What do you type?
  \item Which configuration does a device load on boot, and which one are you
        editing when you type a command?
  \item Why must a hostname and domain name be set before generating an RSA key?
\end{tightnum}
\end{checkpoint}

\begin{chaptersummary}
\begin{tight}
  \item The prompt tells you the mode, and the mode decides which commands are
        legal. \cmd{end} returns to privileged EXEC from anywhere; \cmd{do} runs
        one EXEC command without leaving.
  \item Commands take effect in the running configuration immediately and are
        lost on reload until copied to the startup configuration.
  \item Verify access first, save second. An unsaved risky change is a safety net.
  \item \cmd{enable secret}, not \cmd{enable password}. SSH version 2 with a
        2048-bit key, and \cmd{transport input ssh} to shut off Telnet.
  \item Filtering with \cmd{| include}, \cmd{| section} and \cmd{| begin} is what
        makes device output usable at speed.
\end{tight}
\end{chaptersummary}

\begin{reviewq}
  \item \textbf{(MCQ)} Which prompt indicates global configuration mode?
        (a)~\cmd{SW1>} (b)~\cmd{SW1\#} (c)~\cmd{SW1(config)\#} (d)~\cmd{SW1(config-if)\#}
  \item \textbf{(MCQ)} Which command shows only the lines of the running
        configuration that mention a trunk?
        (a)~\cmd{show run | begin trunk} (b)~\cmd{show run | include trunk}
        (c)~\cmd{show run | section trunk} (d)~\cmd{show run trunk}
  \item \textbf{(Short)} Explain the difference between \cmd{| include} and
        \cmd{| section}, with an example where the answers differ.
  \item \textbf{(Short)} Give the professional argument for \emph{not} saving a
        configuration immediately after making a change to a remote device.
  \item \textbf{(Applied)} Write the complete configuration to take a
        factory-default switch to the point where you can administer it by SSH
        from another subnet, and state which single omission would leave Telnet
        available.
\end{reviewq}
